\documentclass[11pt]{article}
\usepackage[T1]{fontenc}
\usepackage{lmodern}
\usepackage{amsmath,amssymb,amsthm}
\usepackage[margin=1in]{geometry}
\usepackage{microtype}
\usepackage[hidelinks,pdfusetitle]{hyperref}

\newtheorem{theorem}{Theorem}

\newcommand{\Unif}{\operatorname{Unif}}
\newcommand{\E}{\mathbb{E}}

\title{Exact Encoding Ladders for a Uniform Bandwidth Law\\
with an Endpoint Atom}
\author{Xiaojian Zeng\\
PBC School of Finance, Tsinghua University, 100190\\
\texttt{xiaojianzeng@foxmail.com}}
\date{}

\begin{document}
\maketitle

\begin{abstract}
For an encoding ladder \(0=R_0<R_1<\cdots<R_n\leq 1\), a client with
bandwidth \(B\) selects the largest rung not exceeding \(B\), incurring the
one-sided loss \(B-R(B)\).  We solve the finite-\(n\) problem when the bandwidth
law is \((1-\alpha)\Unif[0,1]+\alpha\delta_1\), for every \(n\geq1\) and
\(0\leq\alpha<1\).  The unique optimum changes active set at
\(\alpha=1/(n+1)\): below the transition all \(n+1\) geometric gaps are
positive, while at and above it the top rung is \(1\).  We give the exact
minimum, prove attainment without relaxing the strict ladder inequalities,
and check the low-dimensional and endpoint cases.  This is only the
uniform-plus-single-endpoint-atom specialization of the broader problem for
atomic bandwidth laws.
\end{abstract}

\section{Problem and context}

Fix an integer \(n\geq1\) and \(0\leq\alpha<1\).  Let
\[
 P_\alpha=(1-\alpha)\Unif[0,1]+\alpha\delta_1
\]
and choose \(0=R_0<R_1<\cdots<R_n\leq1\).  At bandwidth \(B\), set
\[
 R(B)=\max\{R_i:R_i\leq B,\ 0\leq i\leq n\}.
\]
The goal is to minimize \(\E[B-R(B)]\), including the equality and
attainment cases under the displayed strict inequalities.

Gray and Neuhoff survey standard scalar quantizers with cell thresholds and
separately chosen reproduction levels~\cite{GrayNeuhoff1998}.  In contrast,
Reznik's adaptive-streaming formulation pins each reproduction value to the
left edge of its cell~\cite{Reznik2026}.  In the uniform,
atom-free, linear-quality case, his Example~1 gives equal cells and loss
\(1/(2(n+1))\); his Remark~1 separately identifies bandwidth laws with atoms
as part of a broader extension problem.  One-way substitution problems of a
related finite character also occur in the assortment literature, where
Sadowski used dynamic programming~\cite{Sadowski1959}.  Here the particular
endpoint mixture makes an exact finite-dimensional convex reduction possible.

An exact sharp-constant calculation that reduces to a low-dimensional convex
minimization, separates endpoint and interior cases, and then classifies
equality also appears in a different matrix setting in Zeng, Liu, and
Ratnavelu~\cite{202608.1346}.
For the present problem the reduction is elementary: Cauchy--Schwarz and a
one-variable quadratic give the entire answer.  No conclusion below concerns
an arbitrary atomic distribution.

\section{Exact solution}

\begin{theorem}
For every \(n\geq1\) and \(0\leq\alpha<1\), the infimum is attained and equals
\[
 L^*_{n,\alpha}=
 \begin{cases}
 \displaystyle
 \frac{1-(n+1)\alpha^2}{2(n+1)(1-\alpha)},
 &0\leq\alpha\leq\dfrac1{n+1},\\[7pt]
 \displaystyle
 \frac{1-\alpha}{2n},
 &\dfrac1{n+1}\leq\alpha<1.
 \end{cases}
\]
The attaining ladder is unique.  For \(1\leq i\leq n\), it is
\[
 R_i=
 \begin{cases}
 \displaystyle\frac{i}{(n+1)(1-\alpha)},
 &0\leq\alpha\leq\dfrac1{n+1},\\[7pt]
 \displaystyle\frac{i}{n},
 &\dfrac1{n+1}\leq\alpha<1.
 \end{cases}
\]
At \(\alpha=1/(n+1)\), the two descriptions coincide.
\end{theorem}

\begin{proof}
Put \(R_{n+1}=1\) and define
\[
 \Delta_i=R_{i+1}-R_i\quad(0\leq i<n),
 \qquad \Delta_n=1-R_n.
\]
The original ladder constraints are equivalent to
\[
 \Delta_i>0\ (0\leq i<n),\qquad
 \Delta_n\geq0,\qquad
 \sum_{i=0}^{n}\Delta_i=1.                                      \tag{1}
\]
Thus \(\Delta_n=0\) is feasible and means \(R_n=1\), whereas a zero
among \(\Delta_0,\ldots,\Delta_{n-1}\) would violate strict increase.

On a uniform draw, the client uses \(R_i\) on
\([R_i,R_{i+1})\).  Hence the contribution of that cell is
\[
 \int_{R_i}^{R_{i+1}}(B-R_i)\,dB=\frac{\Delta_i^2}{2}.
\]
The convention at a cell endpoint is immaterial for the continuous part.
At the atom \(B=1\), however, the client uses \(R_n\), so the loss is exactly
\(1-R_n=\Delta_n\), including when this quantity is zero.  Consequently the
expected loss is
\[
 F(\Delta)=\frac{1-\alpha}{2}\sum_{i=0}^{n}\Delta_i^2
            +\alpha\Delta_n.                                   \tag{2}
\]
This derives both the objective and the feasible set directly from the
original problem.

Let \(t=\Delta_n\).  By (1), \(0\leq t<1\).  At fixed \(t\), the first
\(n\) gaps have sum \(1-t\), and Cauchy--Schwarz gives
\[
 \sum_{i=0}^{n-1}\Delta_i^2\geq\frac{(1-t)^2}{n},                \tag{3}
\]
with equality if and only if
\(
 \Delta_0=\cdots=\Delta_{n-1}=(1-t)/n.
\)
This equality point respects all strict inequalities because \(t<1\).
Thus no closure of the feasible ladder set has been introduced: the original
problem is equivalent to minimizing
\[
 g(t)=\frac{1-\alpha}{2}
      \left(\frac{(1-t)^2}{n}+t^2\right)+\alpha t,
 \qquad 0\leq t<1.                                               \tag{4}
\]
Since
\[
 g''(t)=\frac{(1-\alpha)(n+1)}{n}>0,
\]
this quadratic is strictly convex.  Its unconstrained critical point is
\[
 t_0=\frac{1-(n+1)\alpha}{(n+1)(1-\alpha)}.                      \tag{5}
\]
If \(0\leq\alpha<1/(n+1)\), then \(0<t_0<1\), so \(t_0\) is the unique
minimum.  Equation (3) then makes every first gap equal to
\[
 \frac{1-t_0}{n}=\frac{1}{(n+1)(1-\alpha)}.                      \tag{6}
\]
If \(\alpha\geq1/(n+1)\), then
\[
 g'(0)=\frac{(n+1)\alpha-1}{n}\geq0,
\]
so the unique constrained minimum is \(t=0\); every first gap is then
\(1/n\).  Both candidates satisfy (1), proving attainment in the original
partly open feasible set.  Substitution into (4) gives the two displayed
values.  Summing the gaps gives the displayed rungs.  Finally, strict
convexity of \(g\) and the equality condition in (3) prove uniqueness.
\end{proof}

\section{Direct checks and endpoint behavior}

For \(n=1\), write \(R_1=r\) and \(c=1-\alpha\).  Direct integration in the
original variable gives
\[
 F_1(r)=\frac c2\bigl(r^2+(1-r)^2\bigr)+\alpha(1-r),
 \qquad 0<r\leq1.
\]
Its derivative is \(c(2r-1)-\alpha\).  The stationary point
\(r=1/(2c)\) is feasible exactly for \(\alpha\leq1/2\); thereafter the
unique minimum is \(r=1\).  The corresponding values are
\((1-2\alpha^2)/(4c)\) and \(c/2\), respectively.

For \(n=2\), writing \(R_1=x\) and \(R_2=y\) gives
\[
 F_2(x,y)=\frac c2\bigl(x^2+(y-x)^2+(1-y)^2\bigr)
           +\alpha(1-y).
\]
The interior stationarity equations are \(2x-y=0\) and
\(c(2y-x-1)-\alpha=0\), so
\((x,y)=(1/(3c),2/(3c))\), feasible exactly through
\(\alpha=1/3\).  Above that transition, the active face \(y=1\) has the
unique minimum \(x=1/2\).  The two values are
\((1-3\alpha^2)/(6c)\) and \(c/4\).

At \(\alpha=0\), all \(n+1\) gaps equal \(1/(n+1)\), recovering the
atom-free value \(1/(2(n+1))\).  At the transition, the last gap reaches
zero, both ladder formulas give \(R_i=i/n\), and both value formulas give
\(1/(2(n+1))\).  As \(\alpha\uparrow1\), the optimizer remains
\(R_i=i/n\) while the minimum \((1-\alpha)/(2n)\) tends to zero.  The case
\(\alpha=1\) is excluded from the theorem.  Considered separately, it would
minimize \(1-R_n\), so every strictly increasing ladder with \(R_n=1\) would
attain zero; this explains why uniqueness is lost at the omitted endpoint.

\section*{Scope and assistance disclosure}

The argument determines only the stated uniform-plus-endpoint-atom model.
It neither proves a general theorem for atomic bandwidth distributions nor
establishes a priority claim beyond the literature compared above.

AI assistance was used in proof exploration, a separate mathematical audit,
manuscript drafting, and bibliographic screening.  The author is responsible
for the final text.  Every mathematical step needed for the result is given
analytically above; no numerical optimization, computer algebra, or
computer-generated certificate is used as proof.

\bibliographystyle{plain}
\bibliography{references}

\end{document}
